% s03-3-1-conservativity.tex —— 3.3 引言与 3.3.1 保守性（印刷页 92–97）

\section{Bezrukavnikov 等价}
\label{sec:b}

\subsection{保守性}
\label{sec:conservativity}

本小节说明我们处理因子化 $\Sph^{\spec}_{\check G}$-模范畴的方法。

设 $C_1\to C_2$ 是 $\FactAlgCat$ 中的态射。考虑 \cite{CF} 中定义的因子化限制
（2-）函子
\[
\FactRes_{C_1\to C_2}\colon
C_2\text{-}\FactModCat_{x_0}\longrightarrow C_1\text{-}\FactModCat_{x_0}.
\]
我们将证明，对因子化函子 $\Rep(\check G)\to \Sph^{\spec}_{\check G}$ 而言该函子是保守的。
为此先需要一个引理。

\begin{lemma}\label{lem:3.3.1.1}
设空间 $X$ 使得 $X_{\dR}$ 是 1-仿射的。再设
\[
X=U\sqcup Z,
\]
其中 $U$ 开、$Z$ 闭，并设 $U_{\dR}$ 与 $Z_{\dR}$ 也都是 1-仿射的。记
$\CrysCat(X)=\ShvCat(X_{\dR})$。考虑对应
\begin{equation}
\CrysCat(X)\longrightarrow
\bigl(\CrysCat(U)\times \CrysCat(Z)\bigr)
\times_{\DGCat^{\times 2}}\DGCat^{\Delta^1},
\label{eq:3.3.1}
\end{equation}
\begin{equation}
C\longmapsto (C_U,\,C_Z,\,i^{*}j_{*}).
\label{eq:3.3.2}
\end{equation}
它是保守的。
\end{lemma}

\begin{proof}
用粘贴（recollement）即得。
\end{proof}
\srcpage{92}

\begin{proposition}\label{prop:3.3.1.2}
记
\[
\mathrm{nv}\colon \Rep(\check G)\longrightarrow \Sph^{\spec}_{\check G}
\]
为 \cite[附录 B.2.13]{ABC+24a} 中构造的左伴随。沿 $\mathrm{nv}$ 的因子化限制
\[
\FactRes_{\Rep(\check G)\to \Sph^{\spec}_{\check G}}\colon
\Sph^{\spec}_{\check G}\text{-}\FactModCat_{x_0}\longrightarrow
\Rep(\check G)\text{-}\FactModCat_{x_0}
\]
是保守的。
\end{proposition}

\begin{proof}
设
\[
F\colon M\longrightarrow N
\]
是因子化 $\Sph^{\spec}_{\check G}$-模范畴的态射，使得
\[
\FactRes_{\Rep(\check G)\to \Sph^{\spec}_{\check G}}(F)\colon
\FactRes_{\Rep(\check G)\to \Sph^{\spec}_{\check G}}(M)\longrightarrow
\FactRes_{\Rep(\check G)\to \Sph^{\spec}_{\check G}}(N)
\]
是因子化 $\Rep(\check G)$-模范畴的等价。要证 $F$ 是等价。只需证它作为
$\Ran_{x_0}$ 上的范畴层之间的态射是等价，而这又只需证对每个 $n\in \mathbb{N}$，
\[
F^n_X\colon M^n_X\longrightarrow N^n_X
\]
是等价。

我们对 $n$ 作归纳：$n=0$ 的情形是直接的，因为
\[
F_{x_0}\colon M_{x_0}\longrightarrow N_{x_0}
\]
\srcpage{93}
恰好就是
\[
\FactRes_{\Rep(\check G)\to \Sph^{\spec}_{\check G}}(F)^{x_0}\colon
\FactRes_{\Rep(\check G)\to \Sph^{\spec}_{\check G}}(M)^{x_0}
\xrightarrow{\ \sim\ }
\FactRes_{\Rep(\check G)\to \Sph^{\spec}_{\check G}}(N)^{x_0}.
\]
下面设对一切 $m<n$ 归纳假设成立，来证 $n$ 的情形。记
$U=X\setminus\{x_0\}$，注意
\[
X^n=\bigl(U^n\setminus \Delta_n\bigr)\sqcup
\bigl(\Delta_n\cup (X^n\setminus U^n)\bigr).
\]
由因子化性，
\[
F_{U^n\setminus\Delta_n}\colon
M_{U^n\setminus\Delta_n}\longrightarrow N_{U^n\setminus\Delta_n}
\]
可等同于等价
\[
\mathrm{id}_{\Sph^{\spec,\otimes n}_{\check G}}\otimes F_{x_0}\colon
\Sph^{\spec,\otimes n}_{\check G}\otimes D\text{-}\mathrm{mod}(U^n\setminus\Delta_n)
\otimes M_{x_0}
\xrightarrow{\ \sim\ }
\Sph^{\spec,\otimes n}_{\check G}\otimes
D\text{-}\mathrm{mod}(U^n\setminus\Delta_n)\otimes N_{x_0}.
\]
另一方面，$\Delta_n\cup (X^n\setminus U^n)$ 是超平面 $X^{n-1}$ 的并（且这些超平面到
$\Ran_{x_0}$ 的映射是相容的），故由归纳假设，
\[
F_{\Delta_n\cup(X^n\setminus U^n)}\colon
M_{\Delta_n\cup(X^n\setminus U^n)}\longrightarrow
N_{\Delta_n\cup(X^n\setminus U^n)}
\]
是等价。引用上述引理，只需说明宽松交换图
\[
\begin{tikzcd}[column sep=large,row sep=large]
N_{U^n\setminus\Delta_n}
\arrow[r,"i^{*}j_{*}"]
\arrow[d,"\sim"'] &
N_{\Delta_n\cup(X^n\setminus U^n)}
\arrow[d,"\sim"]\\
M_{U^n\setminus\Delta_n}
\arrow[r,"i^{*}j_{*}"'] &
M_{\Delta_n\cup(X^n\setminus U^n)}
\end{tikzcd}
\]
\srcpage{94}
严格交换。然而这可由相应的立方图推出：该立方图的两个侧面是上述两个方块的
$\FactRes$ 版本，其顶部方块严格交换，而竖直箭头
\[
\FactRes_{\Rep(\check G)\to \Sph^{\spec}_{\check G}}(M)
_{U^n\setminus\Delta_n}\longrightarrow M_{U^n\setminus\Delta_n}
\]
可等同于
\[
\Rep(\check G)^{\otimes n}\otimes D\text{-}\mathrm{mod}(U^n\setminus\Delta_n)
\otimes M_{x_0}\longrightarrow
\Sph^{\spec,\otimes n}_{\check G}\otimes D\text{-}\mathrm{mod}(U^n\setminus\Delta_n)
\otimes M_{x_0},
\]
它生成目标。
\end{proof}

\begin{remark}\label{rem:3.3.1.3}
虽然文献中没有记载，但对任何带连续右伴随的因子化函子 $C\to D$，
$\FactRes$ 都是保守的。由于我们的论证不需要这一点，此处不去证明。
\end{remark}
\srcpage{95}

设
\[
F\colon M_1\longrightarrow M_2
\]
是 $\Sph^{\spec}_{\check G}$-$\FactModCat_{x_0}$ 中的态射。若
$\FactRes_{\Rep(\check G)\to \Sph^{\spec}_{\check G}}(M_1)$ 与
$\FactRes_{\Rep(\check G)\to \Sph^{\spec}_{\check G}}(M_2)$ 作为因子化
$\Rep(\check G)$-模范畴都是可熔的，则要证 $F$ 是等价，只需证
\[
\FactRes_{\Rep(\check G)\to \Sph^{\spec}_{\check G}}(F)^{x_0}=F^{x_0}
\]
是等价。

此外，若 $M$ 是因子化 $\Sph^{\spec}_{\check G}$-模范畴，则存在因子化
$\Sph^{\spec}_{\check G}$-子模范畴
\[
M_{\Sph^{\spec}_{\check G}\text{-}\QLisse}\subset M,
\]
其定义与 $\Rep(\check G)$ 的情形平行，即
\begin{equation}
M_{\Sph^{\spec}_{\check G}\text{-}\QLisse}\big|_{U^n\setminus\Delta_n}
\simeq
\bigl(\Sph^{\spec,\otimes n}_{\check G}\otimes \QLisse(U)^{\otimes n}
\otimes M_{x_0}\bigr)\big|_{U^n\setminus\Delta_n}
\label{eq:3.3.3}
\end{equation}
\begin{equation}
\subset
\bigl(\Sph^{\spec,\otimes n}_{\check G}\otimes D\text{-}\mathrm{mod}(U)^{\otimes n}
\otimes M_{x_0}\bigr)\big|_{U^n\setminus\Delta_n}
\simeq M\big|_{U^n\setminus\Delta_n}.
\label{eq:3.3.4}
\end{equation}
若 $\FactRes_{\Rep(\check G)\to \Sph^{\spec}_{\check G}}(M)$ 可熔，则构造以
$M_{\Sph^{\spec}_{\check G}\text{-}\QLisse}$ 为源的因子化函子，就足以构造以 $M$ 为源的
因子化函子。

下面这一结果将出现在 \cite{BCG26} 中，它重新诠释了目前的讨论。

\begin{proposition}\label{prop:3.3.1.4}
记
\[
\mathrm{Inertia}=Y\times_{Y\times Y}Y
\]
为 $Y=\LS^{\mathrm{restr}}_{\check G}(D^\times_{x_0})$ 的惯性叠。则存在全忠实的
（2-）函子
\[
\mathrm{Fact}^{\mathrm{inertia}}\colon
\QCoh(\mathrm{Inertia})\text{-}\ModCat\longrightarrow
\Sph^{\spec}_{\check G}\text{-}\FactModCat_{x_0},
\]
其本质像恰由那些因子化限制到 $\Rep(\check G)$ 后可熔的因子化
$\Sph^{\spec}_{\check G}$-模范畴组成。
\end{proposition}
\srcpage{96}

\begin{definition}\label{def:aff-spec}
由于有（显然相容的）映射
\[
(\check{\mathfrak{n}}/\check B)_{x_0}\longrightarrow
\LS^{\mathrm{restr}}_{\check G}(D^\times_{x_0}),
\]
以及
\[
\check{\mathfrak{g}}/\check G\longrightarrow
\LS^{\mathrm{restr}}_{\check G}(D^\times_{x_0})\longrightarrow
\LS^{\mathrm{restr}}_{\check G}(D^\times_{x_0})\times
\LS^{\mathrm{restr}}_{\check G}(D^\times_{x_0}),
\]
可以定义因子化 $\Sph^{\spec}_{\check G}$-模范畴 $\Aff^{\spec}_{\check G}$ 为
\[
\Fact^{\mathrm{inertia}}
\Bigl(\IndCoh\bigl((\check{\mathfrak{n}}/\check B)_{x_0}
\times_{\check{\mathfrak{g}}/\check G}
(\check{\mathfrak{n}}/\check B)_{x_0}\bigr)\Bigr).
\]
\end{definition}

\begin{proposition}\label{prop:3.3.1.6}
借助因子化范畴的等价 $\Sat_G$，可以把 $\Aff_G$ 视为因子化
$\Sph^{\spec}_{\check G}$-模范畴。于是
$\FactRes_{\Rep(\check G)\to \Sph^{\spec}_{\check G}}(\Aff^G)$ 可熔。
\end{proposition}

\begin{proof}
与 $\Whit^{!}(\Fl_G)$ 的情形一样，我们再次使用 $D\text{-}\mathrm{mod}(\Fl_G)$ 的可熔性，
但这次是对 Iwahori 取不变量。由构造即得。
\end{proof}

\begin{proposition}\label{prop:3.3.1.7}
$\FactRes_{\Rep(\check G)\to \Sph^{\spec}_{\check G}}
(\Aff^{\spec}_{\check G})$ 可熔。
\end{proposition}
